% Manuscript for corrected canonical run 448f32386dfe45588fefa84c5d486703 (result version m5-seed0-v2).
% Every number below is either a generated macro (paper/generated/results_macros.tex,
% produced by scripts/make_paper_assets.py from the score/protein outputs the paper is built from)
% or a literal. Literals are of three kinds, each labelled where used:
%   (a) values copied from the corrected canonical run's frozen packet
%       (evidence/reviews/m5-canonical-448f32386dfe45588fefa84c5d486703/);
%   (b) historical values of the original unseeded M5 fit (run 29879296d87343898100c190308744a1);
%   (c) historical values from the failed reproduction 78542e9d412548edbfe42b38d452f7a8
%       (evidence/reviews/m5-failure-78542e9d412548edbfe42b38d452f7a8/ and
%       evidence/reviews/m5-reproduction-failure-review.md).
% Core claims: evidence/claims.json; supersession map: evidence/claims-history/m5-seed0-v2-supersession.json;
% deviation record: evidence/protocol-deviations-m5-seed0-v2.md.
% Build: python scripts/build_paper.py (Tectonic, run from paper/).
\documentclass[11pt]{article}

\usepackage[a4paper,margin=2.3cm]{geometry}
\usepackage[T1]{fontenc}
\usepackage{lmodern}
\usepackage{amsmath}
\usepackage{amssymb}
\usepackage{array}
\usepackage{float}
\usepackage{caption}
\usepackage[numbers,sort&compress]{natbib}
\usepackage[hidelinks]{hyperref}
\usepackage{xurl}
\setlength{\emergencystretch}{3em}

% Generated assets (booktabs, graphicx, \ctres, \ctci) and recorded values.
\input{generated/asset_preamble.tex}
\input{generated/results_macros.tex}

\usepackage{adjustbox}
% Generated tabulars are bare; scale each to fit the text block in both directions.
\newcommand{\gentable}[1]{\begin{adjustbox}{max width=\textwidth,max totalheight=0.82\textheight}\input{generated/#1}\end{adjustbox}}
\newcommand{\gd}{$\gamma\delta$~T}
\newcommand{\runid}{\nolinkurl{448f32386dfe45588fefa84c5d486703}}
\newcommand{\origrunid}{\nolinkurl{29879296d87343898100c190308744a1}}
\newcommand{\failrunid}{\nolinkurl{78542e9d412548edbfe42b38d452f7a8}}
\newcommand{\code}[1]{\texttt{#1}}
\captionsetup{font=small,labelfont=bf}

% Reproduction status, set at build time from the generated comparison macro.
% make_paper_assets.py defines \ctres{comparison}{report}{overall} as \textit{not available}
% when no comparison report was supplied (--no-comparison), and as the report's "overall"
% field otherwise. The status text follows that macro and never predicts the outcome.
% An internal tiered comparison is never described as independent verification.
\makeatletter
\def\ct@notavailable{\textit{not available}}
\newif\ifctreport
\expandafter\ifx\csname ctv@comparison@report@overall\endcsname\relax
  \ctreportfalse
\else
  \expandafter\ifx\csname ctv@comparison@report@overall\endcsname\ct@notavailable
    \ctreportfalse
  \else
    \ctreporttrue
  \fi
\fi
\makeatother
\newcommand{\reprostatus}{\ifctreport a seed-controlled reproduction of the corrected M5 results was run and compared with the corrected canonical run by the repository's tiered comparison; that internal comparison reports overall status ``\ctres{comparison}{report}{overall}''. This is an internal tolerance comparison by the same pipeline, not independent verification\else no comparison report for the corrected canonical run was supplied to this build, so the seed-controlled reproduction of the corrected M5 results is pending\fi}
\newcommand{\Reprostatus}{\ifctreport A seed-controlled reproduction of the corrected M5 results was run and compared with the corrected canonical run by the repository's tiered comparison; that internal comparison reports overall status ``\ctres{comparison}{report}{overall}''. This is an internal tolerance comparison by the same pipeline, not independent verification\else No comparison report for the corrected canonical run was supplied to this build, so the seed-controlled reproduction of the corrected M5 results is pending\fi}

\title{Annotation transfer and validation-selected abstention for held-out blood single-cell studies: a recorded benchmark of six matched and two released annotators}
\author{Rewire Bio}
\date{Manuscript draft, result version m5-seed0-v2. Results of record from corrected canonical run \runid. \Reprostatus.}

\begin{document}
\maketitle

\begin{abstract}
\noindent When a new blood or PBMC single-cell study is annotated from a reference, the analyst must decide which annotator to use and which labels to accept. We report a locally frozen protocol that trains six annotators on identical reference cells (the matched track) and runs two released models as distributed (the practical track). All are evaluated on four whole studies held out from the reference. A confidence threshold was fixed on two separate validation studies, at the point that accepted 90\% of validation cells (OP-cov) or kept validation accepted error at or below 5\% (OP-err). Accuracy is agreement with author labels mapped to the Cell Ontology; no labels were adjudicated.

On \ctres{score}{info}{test-donors} test donors and 19{,}018 scorable natural-sample cells, the four learned matched methods had similar closed-set macro-F1 (logistic regression \ctres{A}{M4}{macro-f1-closed}~\ctci{A}{M4}{macro-f1-closed}); the paired intervals for the other three against it all include zero, so no difference was resolved at this precision. At OP-cov, logistic regression accepted \ctres{A}{M4}{coverage-at-OPcov} of test cells with accepted error \ctres{A}{M4}{accepted-error-at-OPcov}. No method reached 0.90 test coverage. Between-study spread was much larger than the pooled accepted-error interval: logistic-regression accepted error ranged from 0.021 to 0.179 across the four studies. After plasmacytoid dendritic cells, antibody-secreting cells and MAIT cells were removed from the reference, between \ctres{B}{M4}{unknown-false-accept-at-OPcov} and \ctres{B}{M3}{unknown-false-accept-at-OPcov} of those cells were still accepted with a known label, usually the nearest relative.

The CellTypist annotator trained on our reference (M5) was originally fitted without a seed, although the protocol specified seed~0. A reproduction attempt found that its per-cell outputs could not be repeated within the pre-specified tolerances. With user approval, both M5 fits were redone with seed~0 and no other change. The M5 values reported here are these corrected, versioned results; the other methods' outputs are reused unchanged from the original recorded run. Intervals are donor-bootstrap intervals within the four fixed studies, conditional on fixed validation thresholds and one training seed, so they say nothing about uncertainty across new studies or seeds. The natural-unknown set contained only 20 erythroid cells and supports counts only. The practical track and the protein and platform checks are descriptive. Reproduction status for this build: \reprostatus.
\end{abstract}

\section{Introduction}

Reference-based annotation assigns each query cell the label of the most similar reference population. Large released models now make this a one-line operation. CellTypist ships immune classifiers trained on cross-tissue atlases \citep{dominguezconde2022celltypist,celltypistmodels}, and scTab ships a classifier trained on the 15.2-million-cell training split of a 22.2-million-cell CELLxGENE corpus \citep{fischer2024sctab}. Most published accuracy figures, however, come from donor-level splits inside one data snapshot. The scTab authors chose donor holdout deliberately, as a compromise between random and study-level splits \citep{fischer2024sctab}. When models trained on the scTab corpus were tested on 21 studies added in a later Census release, macro-F1 fell from 80--84\% to 52--57\% \citep{cultrera2026hce}. A new cohort resembles the second setting more than the first.

The second decision, which labels to accept, has received less systematic attention. A rejection option can fail when a whole population is missing from the reference \citep{abdelaal2019}. Hierarchical and partial rejection has been studied within datasets \citep{theunissen2024}. Conformal methods give guarantees that depend on exchangeability between reference and query \citep{lopezdecastro2025}. Method-intrinsic confidences are calibrated differently across methods \citep{ergen2024popv}, and uncertainty for atlas-level transfer has been measured on a single lung atlas \citep{engelmann2022}. Selective-classification theory provides the vocabulary used here: a classifier plus an acceptance rule, evaluated by coverage and selective risk \citep{geifman2017}.

This paper reports one frozen benchmark designed around these two decisions. It holds out whole studies, fixes thresholds on separate validation studies, and evaluates both matched and released annotators at those fixed operating points. We aim to describe what happened under one stated reference, feature set, default configuration and seed. We do not rank methods or make claims about architectures. Section~\ref{sec:design} states the design, Section~\ref{sec:results} the results, Section~\ref{sec:limits} the limitations, Section~\ref{sec:repro} the reproduction status and commands, and Section~\ref{sec:m5} the M5 seed correction.

\paragraph{Status of the evidence.} The results of record come from corrected canonical run \runid{} (result version \code{m5-seed0-v2}). That run refitted only M5, with seed~0, and reused the completed, hash-verified outputs of every other method from the original recorded run \origrunid. Its comparison manifest records \verb|mode: R| and \verb|fresh_execution: false|, because most steps were adopted rather than executed (Section~\ref{sec:provenance}). The original M5 results, which came from an unseeded fit, are retained as historical records and are reported here only where labelled as such (Section~\ref{sec:m5}). \Reprostatus.

\paragraph{Which numbers are which.} Numbers produced by result macros, the generated tables and the figures are regenerated from whichever score and protein outputs the paper is built from. Numbers typed directly into the text fall into three groups, and each is labelled where it appears.
\begin{itemize}
\item \textbf{Corrected-run literals.} Per-study values, cell and donor counts, per-class recalls and precisions, calibration values quoted in prose and the protein denominators were checked for this revision against the corrected run's \nolinkurl{per_study_test.csv}, \nolinkurl{per_class_test.csv}, \nolinkurl{calibration_test.csv} and \nolinkurl{results.json}. Numbers in figure and table captions (for example 94.3\%/92.4\%, 0.798, 0.847, 15.3\% and 0.99998/0.99992) were checked against the same files.
\item \textbf{Original-run literals for files not in the corrected packet.} The platform table, the per-class unknown-label distributions (most frequent predicted labels for removed and erythroid cells) and the resource-use records come from the original run \origrunid. For M1--M4, M6, P1 and P2 these rest on the same hash-identical predictions and thresholds, but we did not re-check them against corrected-run files. For M5 they describe the \emph{historical unseeded} fit and are marked as such.
\item \textbf{Historical M5 values} from the original unseeded fit and from the failed reproduction \failrunid{} appear only in Section~\ref{sec:m5} and in sentences that say so.
\end{itemize}
Any ratio that combines a literal with a generated value (the ``almost nine times'' comparison in Section~\ref{sec:results}) may not hold exactly after a rebuild from other outputs.

\section{Study design}\label{sec:design}

\subsection{Question and tracks}

The protocol (version~1.0, frozen locally on 2026-10-05 before any validation, test, platform or CITE-seq expression data were downloaded; not externally registered) asks two things: how to choose a reference and annotation workflow for a new blood/PBMC cohort, and when to accept a label. It keeps two tracks separate.

\begin{itemize}
\item \textbf{Matched track.} Every method is trained by us on identical reference cells, labels and genes. Differences between methods are therefore method effects \emph{under this reference, feature set, default settings and training seed}.
\item \textbf{Practical track.} Two released artefacts are used as distributed: CellTypist \nolinkurl{Immune_All_Low} v2 (P1) \citep{celltypistmodels} and the scTab run5 checkpoint (P2) \citep{fischer2024sctab,sctabrepo}. Their training data, label sets and gene spaces differ from each other and from the matched reference, and their labels pass through our frozen label maps. We therefore draw no conclusion about model architecture from any practical-track result, and the matched track supports none beyond the stated configuration.
\end{itemize}

\subsection{Data}

All RNA data come from CELLxGENE Census LTS \texttt{2025-11-08} \citep{czi2025cellxgene,censusreleases}. Every pull is filtered to \verb|is_primary_data == True|, \verb|disease == 'normal'| and \verb|tissue_general == 'blood'|, and raw counts are used.

\begin{itemize}
\item \textbf{Reference (six studies):} Hao 2021 \citep{hao2021seurat}, OneK1K \citep{yazar2022onek1k}, Perez 2022 \citep{perez2022lupus}, van der Wijst 2021 \citep{vanderwijst2021}, COMBAT \citep{combat2022} and the CVID atlas \citep{rodriguezubreva2022cvid}.
\item \textbf{Validation, used only to set thresholds and M4's regularisation:} the kidney-cancer cohort blood samples \citep{li2022kidney} and a clonal-haematopoiesis study \citep{heimlich2024ch}.
\item \textbf{Test, each held out as a whole study:} the Human Immune Health Atlas (HIHA) \citep{gong2025hiha}, a rheumatoid-arthritis PBMC study (RA) \citep{binvignat2024ra}, a glaucoma PBMC atlas (Glaucoma) \citep{glaucomacollection} and a juvenile-dermatomyositis CITE-seq study (JDM) \citep{rabadam2024jdm}.
\item \textbf{Platform panel (descriptive):} one dataset per assay from a single-donor technology comparison \citep{desimone2025platforms}.
\end{itemize}

\noindent The protocol records that the validation, test and platform studies were absent from scTab's 249 training datasets, share no donor IDs with scTab's donor lookup, and were published after the CellTypist model build (2022-07-16).

Sampling used seed 20261005. For validation and test, the protocol took donors with at least 300 cells, up to 12 per study. It drew a \emph{natural stratum} of 600 cells per donor from all of that donor's cells, preserving composition, and a \emph{rare top-up stratum} of up to 40 further mapped cells per donor and rare class. Unless stated otherwise, every test metric uses the natural stratum only.

The corrected run scored \ctres{score}{info}{test-donors} test donors in total, on the same hash-verified data build (D05) as the original run. A read-only metadata audit of the D05 query files made for the original run (\nolinkurl{evidence/results-29879296/donor-count-check.json}; no rescoring) gives the per-study donor counts: HIHA 12, RA 12, Glaucoma 3 and JDM 5, all with natural-stratum cells, which sum to 32. Their natural strata hold 7{,}200, 7{,}200, 1{,}800 and 3{,}000 cells (600 per donor); fewer are scorable after label mapping. Glaucoma and JDM therefore rest on three and five donors, so their within-study bootstrap has very few clusters. The realised cell counts per study, stratum and label status are in \nolinkurl{cell_counts.csv} of the scoring output and in the \nolinkurl{cell_counts} block of the corrected run's \nolinkurl{results.json}.

\subsection{Labels, classes and arms}

Author \nolinkurl{cell_type_ontology_term_id} labels were mapped to target classes with a frozen table built on the Cell Ontology \citep{diehl2016cl}. Only cells with status \texttt{mapped} are used for training or fine-level scoring.

\paragraph{Arm A (full reference).} The known classes $K$ are target classes with at least 100 reference cells from at least two reference studies. There are 14: ASC, B, CD14 mono, CD16 mono, CD4~T, CD8~T, HSPC, ILC, MAIT, NK, cDC, \gd, pDC and platelet/MK. Closed-set macro-F1 averages over the classes with at least 20 natural test cells. In this run that left 11 classes; HSPC (17 natural cells), ILC (17) and platelet/MK (16) fell below the cut-off. Mapped validation or test cells whose class is outside $K$ are \emph{natural unknowns}.

\paragraph{Arm B (removed types).} M1--M6 are retrained after every pDC, ASC and MAIT cell is removed from the reference, and thresholds are re-tuned on validation cells of the remaining classes. Test cells of the three removed classes, from both strata, are \emph{simulated unknowns}. Arm~B macro-F1 averages over 8 classes and is not comparable with Arm~A macro-F1.

\subsection{Annotators}

All matched methods use the same features $F$, chosen on the reference only: 2{,}000 highly variable genes plus the top 10 marker genes per reference class.
\begin{itemize}
\item \textbf{M1, marker rule.} Mean z-scored expression of 10 markers per class; confidence is top minus second score.
\item \textbf{M2, nearest centroid} on 50 reference PCs with cosine similarity; confidence is top minus second similarity.
\item \textbf{M3, kNN} ($k=15$, cosine); confidence is the vote share of the predicted class.
\item \textbf{M4, multinomial logistic regression,} with $C\in\{0.01,0.1,1\}$ chosen by validation closed-set macro-F1 (the protocol's only tuning budget); confidence is maximum probability.
\item \textbf{M5, CellTypist~1.7.1} trained on our reference \citep{dominguezconde2022celltypist}; confidence is maximum one-vs-rest sigmoid probability. With more than 50{,}000 reference cells in both arms, CellTypist selects scikit-learn's SAG solver; \code{max\_iter} stays at 500 and \code{n\_jobs} at 4. In the results of record the solver receives \code{random\_state=0}; the original recorded fits received none (Section~\ref{sec:m5}).
\item \textbf{M6, scANVI} (scvi-tools 1.4.1) \citep{xu2021scanvi}, with per-study query adaptation on \emph{unlabelled} query cells in the style of scArches \citep{lotfollahi2022scarches}; confidence is maximum probability.
\end{itemize}
P1 and P2 receive inputs in their documented form. Released labels are mapped as \texttt{mapped} (scored at the target level), \texttt{coarser} (never counted as accepted at the target level) or \texttt{outside} (a fine-level error if accepted). Pan-Human Azimuth \citep{sarkar2026panhuman} and scGPT \citep{cui2024scgpt} were not run, for the reasons recorded in the protocol's access-and-feasibility notes. The protocol specifies training seed~0. In the results of record, M2/M3 (PCA), M4, M5 and M6 are seeded with 0, and M1 has no random component. Each method was fitted once per arm for the results of record. The original M5 fits passed no seed, which deviated from the protocol; this is logged in \nolinkurl{evidence/protocol-deviations-m5-seed0-v2.md}.

\subsection{Operating points selected on validation}\label{sec:ops}

For each method and arm, thresholds were chosen on validation cells that are mapped and in $K$. As implemented, these cells include both strata (\ctres{A}{M4}{val-validation-cells} cells in Arm~A, \ctres{B}{M4}{val-validation-cells} in Arm~B). The validation mix is therefore enriched for rare classes relative to the test natural stratum.

\begin{itemize}
\item \textbf{OP-cov:} the threshold $\tau_\mathrm{cov}$ that accepted 90\% of these validation cells. If the method's maximum eligible validation coverage is below 90\%, the protocol accepts every fine-level prediction and reports the shortfall.
\item \textbf{OP-err:} the lowest threshold whose validation accepted error is at most 5\%; if none exists, it is reported as ``not attainable''.
\end{itemize}

\noindent Three implementation facts shape how these operating points should be read.
\begin{enumerate}
\item \textbf{P2 has no OP-cov threshold.} Only a fraction \ctres{practical}{P2}{val-eligible-cap-validation} of P2's validation predictions were fine-level labels eligible for acceptance (its eligible validation cap), which is below 0.90. By the shortfall rule, $\tau_\mathrm{cov}=-\infty$: P2 at ``OP-cov'' accepts every fine-level prediction and abstains on nothing. (The CSV records this as \texttt{-inf}; \texttt{results.json} serialises it as \texttt{null}.)
\item \textbf{M3's vote-share confidence takes only 16 values}, and the threshold accepts all cells with confidence $\ge\tau$. Whole tie blocks are therefore accepted together. M3's validation coverage at $\tau_\mathrm{cov}$ was \ctres{A}{M3}{val-validation-coverage-at-tau-cov} in Arm~A and \ctres{B}{M3}{val-validation-coverage-at-tau-cov} in Arm~B instead of 0.900.
\item \textbf{The bootstrap does not re-estimate thresholds.} They are fixed numbers carried to the test set (Section~\ref{sec:unc}).
\end{enumerate}

\subsection{Metrics and denominators}

All test metrics use scorable cells: mapped, in $K$, natural stratum. Arm~A has 19{,}018 such cells and Arm~B 18{,}628.
\begin{itemize}
\item \textbf{Coverage} is accepted cells divided by scorable cells.
\item \textbf{Accepted error} is wrong accepted cells divided by accepted cells.
\item \textbf{Cross-lineage rate} is accepted errors that cross a fixed lineage group, divided by accepted cells. \textbf{Cross-lineage share} is the same errors divided by all accepted errors.
\item \textbf{Unknown false acceptance} is accepted unknown cells divided by unknown cells. For Arm~A, only natural-stratum unknowns are counted (the approved D-1a scoring correction); for Arm~B, all strata are counted.
\item \textbf{Per-class recall} uses natural and top-up cells; \textbf{precision} uses natural cells only.
\item \textbf{Calibration} is reported as the multiclass Brier score and the 15-bin ECE of the maximum probability \citep{guo2017}.
\end{itemize}

\noindent The risk--coverage curve sorts cells by confidence. In this implementation, AURC is the mean risk multiplied by the final coverage, so the area stops at each method's maximum coverage.

For P2, ``unassigned'' cells at OP-cov are its \texttt{coarser} labels (\ctres{practical}{P2}{coarser-fraction} of test cells), since nothing is abstained.

\subsection{Uncertainty}\label{sec:unc}

Intervals are 95\% percentile intervals from 1{,}000 cluster-bootstrap replicates. Donors are resampled with replacement \emph{within} each of the four test studies, so study membership is fixed. Method differences are paired within replicates against M4, which the frozen protocol named as the comparator before any result existed.

Throughout this paper, every bracketed interval is a \emph{95\% donor-bootstrap percentile interval within the four test studies (\ctres{score}{info}{reps} replicates), conditional on fixed validation thresholds and a single training seed}. The intervals do not include:
\begin{itemize}
\item uncertainty from estimating the thresholds;
\item training-seed variance, which was not estimated for any method (Section~\ref{sec:m5} explains why the M5 history does not provide such an estimate);
\item variation across studies, since only four test studies exist and no study-level interval can be computed.
\end{itemize}

Pooled metrics weight studies by their number of scorable natural cells. In Arm~A these are HIHA 7{,}042, RA 7{,}200, JDM 3{,}000 and Glaucoma 1{,}776, so HIHA and RA together contribute 75\%. Per-study values are therefore given next to pooled values (Section~\ref{sec:results}; full per-study table in \nolinkurl{per_study_test.csv} of the scoring output). No sensitivity analysis excluding HIHA, or reweighting studies equally, was performed; only the per-study values are available.

Paired differences against M4 are reported as the \emph{mean of the paired bootstrap differences} (replicate value minus the replicate M4 value), with the percentile interval of those differences. For the differences quoted in the text, this mean differs from the difference between point estimates by at most 0.001.

\paragraph{Absent classes in bootstrap replicates.} Because donors are resampled, some replicates contain no cell of a rare class or no unknown cell. The two cases are handled differently by the scoring code.
\begin{itemize}
\item \textbf{Macro-F1.} The macro-F1 class set is fixed once from the full sample. A replicate that contains no cell of one of these classes scores that class's F1 as 0 and is kept. This biases replicate macro-F1 downward; how often it happened was not recorded. Macro-F1 intervals are therefore \emph{not} conditional on the class being present, and no replicate was undefined for macro-F1.
\item \textbf{Undefined replicates.} Table~\ref{tab:nan} counts the replicates in which a metric was undefined. They arose only in two ways: Arm~A and practical-track natural-unknown false acceptance in the 126 replicates that contained no unknown cell, and OP-err metrics for M3 and M6 (1{,}000 replicates), where the operating point was not attainable. No per-class metric was bootstrapped.
\end{itemize}
The natural-unknown interval therefore rests on 874 replicates and is conditional on an unknown cell being present. For example, M4's value is 0.15 with interval [0.105, 1.000]: the interval contains the point estimate but spans almost the whole range. A percentile interval built this way could in principle exclude its point estimate, but none of the natural-unknown intervals in the corrected run did. We treat these intervals as uninformative.

\subsection{Secondary analyses}

\paragraph{Platform.} We report coverage and accepted error at OP-cov for each platform dataset. The panel comes from one donor, so these results are descriptive only.

\paragraph{Protein check.} Accepted RNA predictions at OP-cov are compared with a coarse protein class. The protein class is assigned by six fixed ADT gates (CD4~T, CD8~T, B, NK, CD14 mono, CD16 mono) after per-cell CLR normalisation and a two-component Gaussian mixture per antibody \citep{stoeckius2017citeseq}. The pre-specified 10x files could not be used, and this check ran on replacement inputs under an approved amendment (Section~\ref{sec:amend}).

\subsection{Run provenance}\label{sec:provenance}

The corrected canonical run \runid{} (code revision \nolinkurl{036177f41adb2b2a535e2c4700868ad5807252e3}, completed 2026-10-07) assembled its inputs as follows. Every adopted file is listed with its sha256 in the run's generation receipt and was hash-checked by the repair driver.
\begin{itemize}
\item \textbf{Adopted unchanged from the original run's sources:} the data build D05, the released-model predictions for P1 and P2, the CITE build, all Arm~A and Arm~B fits and predictions for M1--M4 and M6, and the Arm~A CITE predictions for M1--M4 and M6. As in the original run, Arm~A M1--M6 and Arm~B M1--M3 come from the historical matched runs, and Arm~B M4 and M6 come from the earlier recovery run \nolinkurl{6ef427915cfa4087a9bdf8646a0c26f7}. No M5 file from any earlier run was adopted.
\item \textbf{Arm~A M5 fit, continued without re-execution.} The seed-0 Arm~A fit was executed once, in the failed canonical attempt \nolinkurl{3f18e178b78d4d75b595174cfa00cf99}. The corrected run carried that completed fit forward: its receipt records \verb|executed: false|, the source run, the source step hash \nolinkurl{7e52e4a650571bf4c0a82aa870db2efdcad1ffac9d8103a518825ebc626e802f} and the model hash \nolinkurl{111e03c33dc7f8062fbb5106af155720852e2b95a0a067ad08c80836977ef11b}. We do not describe this fit as newly executed in run \runid.
\item \textbf{Executed in run \runid:} the Arm~B M5 fit (seed~0; 501~s; model hash \nolinkurl{15829b4c79f3c1ca510bcc2a7487501d25ca4656e1a88dcbf2916c87adbb2770}), the Arm~A and Arm~B M5 predictions, the Arm~A M5 CITE predictions, scoring of all methods (613~s) and the protein check (11~s).
\end{itemize}

\noindent The receipt records \code{random\_state = 0} and \code{solver = sag} for both corrected M5 models, and gene, class and scaler hashes equal to those of the historical M5 models in each arm. After scoring, the driver's non-M5 equality check compared every non-M5 row of the summary and threshold tables, the cell counts and the bootstrap weights with the original run and recorded \verb|passed|. The sha256 values of \nolinkurl{score_all.py}, \nolinkurl{evaluate.py} and \nolinkurl{protein_check.py} are unchanged from the original run (Section~\ref{sec:repro}). The only scientific source change is the M5 seed argument in \nolinkurl{methods.py} (Section~\ref{sec:m5}).

The original run \origrunid{} was itself a recovery run (``Mode R''): it reused cached and hash-verified data, Arm~A and Arm~B M1--M3 outputs, adopted Arm~B M4--M6 from run \nolinkurl{6ef427915cfa4087a9bdf8646a0c26f7}, and recomputed the CITE build, the Arm~A CITE predictions, scoring and the protein check. Its assembly check (\nolinkurl{evidence/results-29879296/assembly-check.json}) concerns consistency of the assembled inputs and is not an independent reproduction.

\section{Results}\label{sec:results}

\subsection{Closed-set agreement with author labels}

Table~\ref{tab:closed} gives closed-set macro-F1, computed with no abstention. In Arm~A, the four learned matched methods were within 0.005 of one another:
\begin{itemize}
\item M4: \ctres{A}{M4}{macro-f1-closed}~\ctci{A}{M4}{macro-f1-closed};
\item M3: \ctres{A}{M3}{macro-f1-closed};
\item M6: \ctres{A}{M6}{macro-f1-closed};
\item M5 (seed~0): \ctres{A}{M5}{macro-f1-closed}.
\end{itemize}
Their paired differences from M4 all include zero: M3 \ctres{A}{M3}{diff-vs-A-M4-macro-f1-closed}~\ctci{A}{M3}{diff-vs-A-M4-macro-f1-closed}, M5 \ctres{A}{M5}{diff-vs-A-M4-macro-f1-closed}~\ctci{A}{M5}{diff-vs-A-M4-macro-f1-closed} and M6 \ctres{A}{M6}{diff-vs-A-M4-macro-f1-closed}~\ctci{A}{M6}{diff-vs-A-M4-macro-f1-closed}. The M5 interval also included zero for both earlier unseeded M5 fits (Section~\ref{sec:m5}). The centroid and marker baselines had lower agreement: M2 \ctres{A}{M2}{macro-f1-closed} (difference \ctres{A}{M2}{diff-vs-A-M4-macro-f1-closed}~\ctci{A}{M2}{diff-vs-A-M4-macro-f1-closed}) and M1 \ctres{A}{M1}{macro-f1-closed} (difference \ctres{A}{M1}{diff-vs-A-M4-macro-f1-closed}~\ctci{A}{M1}{diff-vs-A-M4-macro-f1-closed}). Under this configuration, no closed-set difference between logistic regression, kNN, CellTypist-trained-on-our-reference and scANVI was resolved at the precision of these intervals. This does not show that the methods are equivalent; no equivalence margin was pre-specified or tested.

The practical-track models were run as distributed, with labels mapped by our frozen tables. P1 had lower closed-set macro-F1 than matched M4, with a paired difference of \ctres{practical}{P1}{diff-vs-A-M4-macro-f1-closed}~\ctci{practical}{P1}{diff-vs-A-M4-macro-f1-closed}. For P2 the paired interval against M4 includes zero, so no difference was resolved: \ctres{practical}{P2}{diff-vs-A-M4-macro-f1-closed}~\ctci{practical}{P2}{diff-vs-A-M4-macro-f1-closed}. Some practical deficits are label-map effects. P1's closed-set recall for \gd{} cells was 0.015, while only \ctres{practical}{P1}{coarser-fraction} of P1's predictions were \texttt{coarser}. P1 therefore assigned most \gd{} cells to other classes or to labels outside $K$ (P1's \texttt{outside} fraction is listed in Table~\ref{tab:practical}), which is consistent with how CellTypist's T-cell labels map to our target classes; we did not test this further. Table~\ref{tab:practical} (Appendix) lists the \texttt{coarser} and \texttt{outside} fractions.

\begin{table}[t]
\centering
\caption{Closed-set macro-F1 (no abstention), AURC, maximum coverage and cross-lineage error at OP-cov. Arm~A macro-F1 averages over 11 classes and Arm~B over 8, so the two arms cannot be compared on this metric. M5 rows are the corrected seed-0 fits. Intervals: 95\% donor-bootstrap percentile intervals within the four test studies, conditional on fixed thresholds and one training seed. AURC stops at each method's maximum coverage, so P2's AURC (maximum coverage 0.847) is not comparable with the other rows, and M3's AURC depends on row order within confidence ties (Section~\ref{sec:rc}).}
\label{tab:closed}
\gentable{table_closed_set.tex}
\end{table}

\paragraph{Per-study heterogeneity.} Pooled values hide large differences between studies (full table: \nolinkurl{per_study_test.csv}). For M4 at OP-cov (point estimates from the corrected run's per-study table; M4 outputs are unchanged from the original run; no per-study interval):

\begin{center}
\small
\begin{tabular}{lrrr}
\toprule
Study & Accepted error & Coverage & Macro-F1 (closed) \\
\midrule
HIHA & 0.021 & 0.940 & 0.942 \\
JDM & 0.056 & 0.906 & 0.886 \\
Glaucoma & 0.144 & 0.847 & 0.707 \\
RA & 0.179 & 0.651 & 0.518 \\
\bottomrule
\end{tabular}
\end{center}

\noindent The pooled M4 accepted error is \ctres{A}{M4}{accepted-error-at-OPcov}~\ctci{A}{M4}{accepted-error-at-OPcov}. The study-to-study range (0.157, a literal) was almost nine times the width of that interval; this ratio mixes a literal with a generated interval and may not hold exactly after a rebuild from other outputs. On closed-set macro-F1, every method in both arms and both tracks followed the same ordering at the extremes: HIHA was highest and RA lowest (for example, among the Arm~A matched methods, RA macro-F1 ranged from 0.353 for M1 to 0.518 for M4; corrected M5 was 0.507). We can say only that agreement with RA author labels was lower. Differences in granularity or definitions are plausible explanations, but we did not test them.

The HIHA result needs a separate caution. The HIHA authors report that their expert annotation was guided by CellTypist \nolinkurl{Immune_All} models and Seurat label transfer \citep{gong2025hiha}. Agreement with HIHA labels therefore partly measures similarity to that labelling pipeline. P1 is one of the guiding models. Yet P1's HIHA macro-F1 (0.842) was below M4's (0.942), so these numbers show no obvious advantage for P1 on HIHA. This does not rule out circularity acting through the Seurat step or through label granularity.

\subsection{Abstention at the validation-selected OP-cov threshold}

Table~\ref{tab:ops} and Figure~\ref{fig:ops} give coverage and accepted error at both operating points. OP-cov is ``the threshold that accepted 90\% of validation cells''. On the test studies, Arm~A coverage at this threshold ranged from \ctres{A}{M5}{coverage-at-OPcov} (M5) to \ctres{A}{M3}{coverage-at-OPcov} (M3), and no method or arm reached 0.900. Validation error at the threshold was higher than test error; for M4 it was \ctres{A}{M4}{val-validation-error-at-tau-cov} on validation against \ctres{A}{M4}{accepted-error-at-OPcov} on test. Both shifts are consistent with the validation mix being enriched for hard, rare classes.

\begin{itemize}
\item \textbf{M4} accepted \ctres{A}{M4}{coverage-at-OPcov}~\ctci{A}{M4}{coverage-at-OPcov} of test cells with accepted error \ctres{A}{M4}{accepted-error-at-OPcov}~\ctci{A}{M4}{accepted-error-at-OPcov}.
\item \textbf{M5 (seed~0)} accepted \ctres{A}{M5}{coverage-at-OPcov} with error \ctres{A}{M5}{accepted-error-at-OPcov}. Its paired coverage difference from M4 was \ctres{A}{M5}{diff-vs-A-M4-coverage-at-OPcov} \ctci{A}{M5}{diff-vs-A-M4-coverage-at-OPcov}. Its paired error difference was \ctres{A}{M5}{diff-vs-A-M4-accepted-error-at-OPcov} \ctci{A}{M5}{diff-vs-A-M4-accepted-error-at-OPcov}.
\item \textbf{M6} accepted more cells (\ctres{A}{M6}{coverage-at-OPcov}) with higher error (\ctres{A}{M6}{accepted-error-at-OPcov}).
\item \textbf{M3} accepted \ctres{A}{M3}{coverage-at-OPcov} with error \ctres{A}{M3}{accepted-error-at-OPcov}. Because of ties, M3's threshold accepted \ctres{A}{M3}{val-validation-coverage-at-tau-cov} of validation cells against the nominal 0.900. Its paired error difference from M4 (\ctres{A}{M3}{diff-vs-A-M4-accepted-error-at-OPcov}~\ctci{A}{M3}{diff-vs-A-M4-accepted-error-at-OPcov}) therefore compares two different coverages.
\end{itemize}

\noindent The M4 and M5 thresholds produced the lowest accepted errors, but at lower coverage than M3 and M6. These are coverage--error trade-offs at fixed thresholds and give no single winner.

For the practical track:
\begin{itemize}
\item \textbf{P1} accepted \ctres{practical}{P1}{coverage-at-OPcov} with error \ctres{practical}{P1}{accepted-error-at-OPcov}.
\item \textbf{P2} accepted \ctres{practical}{P2}{coverage-at-OPcov} with error \ctres{practical}{P2}{accepted-error-at-OPcov}. These describe P2 \emph{with no abstention}: the ``unassigned'' part is its \ctres{practical}{P2}{coarser-fraction} of \texttt{coarser} labels, and no prediction was rejected. The paired P2$-$M4 coverage difference (\ctres{practical}{P2}{diff-vs-A-M4-coverage-at-OPcov}) therefore compares an unthresholded model with a thresholded one.
\end{itemize}

\begin{table}[t]
\centering
\caption{Coverage and accepted error on the test natural stratum at the two validation-selected operating points. M5 rows are the corrected seed-0 fits. ``Not attainable'' means no validation threshold reached accepted error $\le$5\%. P2 has no OP-cov threshold, so its OP-cov row accepts every fine-level label. M3's OP-cov threshold accepted 94.3\% (Arm~A) and 92.4\% (Arm~B) of validation cells because of ties. The OP-err thresholds of M4 and P2 (0.99998 and 0.99992) and of corrected M5 (0.9999998 in both arms) sit at the probability ceiling. Intervals as in Table~\ref{tab:closed}.}
\label{tab:ops}
\gentable{table_operating_points.tex}
\end{table}

\begin{figure}[t]
\centering
\includegraphics[width=0.95\textwidth]{figures/fig_operating_points.pdf}
\caption{Coverage against accepted error at the validation-selected operating points, drawn from the values and intervals of the outputs the paper is built from. Practical-track markers are hollow. The P2 point at OP-cov has no threshold (all fine-level labels accepted). The M3 OP-cov point accepted more than 90\% of validation cells because of tied confidences. Intervals are within-study donor-bootstrap intervals and do not describe variation across new studies.}
\label{fig:ops}
\end{figure}

\subsection{The stricter OP-err threshold transferred unevenly}

OP-err aimed for accepted error of at most 5\% on validation. On the test studies:
\begin{itemize}
\item \textbf{M2} kept \ctres{A}{M2}{coverage-at-OPerr} of cells at \ctres{A}{M2}{accepted-error-at-OPerr} error.
\item \textbf{P1} kept \ctres{practical}{P1}{coverage-at-OPerr} at \ctres{practical}{P1}{accepted-error-at-OPerr}.
\item \textbf{M5} kept \ctres{A}{M5}{coverage-at-OPerr} at \ctres{A}{M5}{accepted-error-at-OPerr}.
\item \textbf{M4} kept \ctres{A}{M4}{coverage-at-OPerr}, with no accepted errors (point and interval \ctci{A}{M4}{accepted-error-at-OPerr}).
\item \textbf{M1} missed the target, at \ctres{A}{M1}{accepted-error-at-OPerr}~\ctci{A}{M1}{accepted-error-at-OPerr} error.
\item \textbf{P2} kept \ctres{practical}{P2}{coverage-at-OPerr} of cells, about 47 cells. It accepted none of the RA cells, so its RA accepted error is undefined.
\item \textbf{M3 and M6:} OP-err was not attainable in either arm.
\end{itemize}
The OP-err thresholds of M4 and P2 (0.99998 and 0.99992) and of corrected M5 (0.9999998 in both arms, from \nolinkurl{thresholds_validation.csv}) sit at the probability ceiling, so their OP-err results are close to degenerate.

The paired M2$-$M4 coverage difference at OP-err, \ctres{A}{M2}{diff-vs-A-M4-coverage-at-OPerr}~\ctci{A}{M2}{diff-vs-A-M4-coverage-at-OPerr}, is a measured difference at these fixed thresholds. It reflects how each confidence score ranked validation errors relative to test errors. We do not read it as evidence that centroid classification is generally safer.

\subsection{Risk--coverage}\label{sec:rc}

Figure~\ref{fig:rc} shows the test risk--coverage curves. Risk at a common coverage of 0.80 can be compared across both tracks, because every method reaches it:

\begin{center}
\small
\begin{tabular}{lcccccccc}
\toprule
& M4 & M6 & M5 & M3$^\dagger$ & P1 & P2 & M2 & M1 \\
\midrule
risk@0.80 & \ctres{A}{M4}{risk-at-0.80} & \ctres{A}{M6}{risk-at-0.80} & \ctres{A}{M5}{risk-at-0.80} & \ctres{A}{M3}{risk-at-0.80} & \ctres{practical}{P1}{risk-at-0.80} & \ctres{practical}{P2}{risk-at-0.80} & \ctres{A}{M2}{risk-at-0.80} & \ctres{A}{M1}{risk-at-0.80} \\
\bottomrule
\end{tabular}
\end{center}

\noindent These values have no interval. $^\dagger$The risk--coverage computation sorts cells with a stable sort and does not average within confidence ties. Inside a tie block the order follows the input row order, which is the order in which studies were concatenated. M3's curve, AURC (\ctres{A}{M3}{aurc}) and risk@$c$ values therefore contain an arbitrary component, and we do not rank M3 on them. Methods with exactly tied probabilities, such as saturated values of 1.0, could be affected in the same way; we did not check this.

P2 cannot reach coverage of 0.90 or more, so its risk at those coverages is ``not reachable''. P1 cannot reach 1.00 (maximum coverage \ctres{practical}{P1}{max-coverage}). Because AURC stops at maximum coverage, P2's AURC (\ctres{practical}{P2}{aurc}) covers only coverage up to \ctres{practical}{P2}{max-coverage}. It cannot be compared with M4's (\ctres{A}{M4}{aurc}) or M6's (\ctres{A}{M6}{aurc}).

\begin{figure}[t]
\centering
\includegraphics[width=0.95\textwidth]{figures/fig_risk_coverage.pdf}
\caption{Risk (accepted error) against coverage on the test natural stratum, from the curves of the outputs the paper is built from. Practical curves are dashed. The P2 curve ends at coverage 0.847 because 15.3\% of its labels are \texttt{coarser}. Ties are not averaged, so the shape of M3's curve within its 16 tie blocks depends on row order.}
\label{fig:rc}
\end{figure}

\subsection{Removed cell types (Arm B)}

Arm~B tested what happens to three specific cell types once they are missing from the reference. There were 2{,}310 test cells of the removed classes: ASC 815, MAIT 748 and pDC 747. Only 390 of them came from the natural stratum; the other 83\% came from the rare top-up stratum. Pooled false acceptance therefore weights the three classes according to the sampling design (about one third each); their natural prevalence plays almost no part.

\begin{table}[t]
\centering
\small
\caption{Arm~B: acceptance of cells whose class (ASC, MAIT or pDC) was removed from the reference (2{,}310 cells, all strata). False acceptance is the fraction of these cells accepted with a known label. AUROC measures how well confidence separates known from removed-type cells; 0.5 is no separation. M5 metric columns are the corrected seed-0 fit. The most frequent predicted labels are from the per-class unknown table of the original run (fraction of that class's cells); for M1--M4 and M6 they rest on the same predictions as the results of record. $^h$For M5 they describe the historical unseeded fit; the corresponding table for the corrected fit was not in the packet supplied for this revision. Intervals as in Table~\ref{tab:closed}.}
\label{tab:armb}
\setlength{\tabcolsep}{4pt}
\begin{tabular}{@{}l>{\raggedright\arraybackslash}p{2.7cm}>{\raggedright\arraybackslash}p{1.7cm}>{\raggedright\arraybackslash}p{1.3cm}>{\raggedright\arraybackslash}p{5.6cm}@{}}
\toprule
Method & False accept.\ at OP-cov & False accept.\ at OP-err & AUROC & Most frequent label (fraction of class): ASC / MAIT / pDC \\
\midrule
M1 & \ctres{B}{M1}{unknown-false-accept-at-OPcov}\newline{\scriptsize\ctci{B}{M1}{unknown-false-accept-at-OPcov}} & \ctres{B}{M1}{unknown-false-accept-at-OPerr} & \ctres{B}{M1}{unknown-auroc} & ILC .269 / CD8~T .305 / cDC .767 \\
M2 & \ctres{B}{M2}{unknown-false-accept-at-OPcov}\newline{\scriptsize\ctci{B}{M2}{unknown-false-accept-at-OPcov}} & \ctres{B}{M2}{unknown-false-accept-at-OPerr} & \ctres{B}{M2}{unknown-auroc} & B .810 / \gd{} .929 / B .451 \\
M3 & \ctres{B}{M3}{unknown-false-accept-at-OPcov}\newline{\scriptsize\ctci{B}{M3}{unknown-false-accept-at-OPcov}} & \ctres{B}{M3}{unknown-false-accept-at-OPerr} & \ctres{B}{M3}{unknown-auroc} & B .866 / \gd{} .914 / cDC .854 \\
M4 & \ctres{B}{M4}{unknown-false-accept-at-OPcov}\newline{\scriptsize\ctci{B}{M4}{unknown-false-accept-at-OPcov}} & \ctres{B}{M4}{unknown-false-accept-at-OPerr} & \ctres{B}{M4}{unknown-auroc} & B .899 / \gd{} .639 / cDC .712 \\
M5 & \ctres{B}{M5}{unknown-false-accept-at-OPcov}\newline{\scriptsize\ctci{B}{M5}{unknown-false-accept-at-OPcov}} & \ctres{B}{M5}{unknown-false-accept-at-OPerr} & \ctres{B}{M5}{unknown-auroc} & $^h$B .852 / \gd{} .595 / cDC .497 \\
M6 & \ctres{B}{M6}{unknown-false-accept-at-OPcov}\newline{\scriptsize\ctci{B}{M6}{unknown-false-accept-at-OPcov}} & \ctres{B}{M6}{unknown-false-accept-at-OPerr} & \ctres{B}{M6}{unknown-auroc} & B .639 / \gd{} .797 / HSPC .959 \\
\bottomrule
\end{tabular}
\end{table}

At the OP-cov threshold, every matched method accepted most cells of the removed types with a known label (Table~\ref{tab:armb}). The fraction ranged from \ctres{B}{M4}{unknown-false-accept-at-OPcov}~\ctci{B}{M4}{unknown-false-accept-at-OPcov} for M4 to \ctres{B}{M3}{unknown-false-accept-at-OPcov} for M3 and \ctres{B}{M6}{unknown-false-accept-at-OPcov} for M6. The most frequent predicted labels (over all cells of each removed class, not only accepted cells; original-run table, historical for M5) were mostly the nearest known relative:
\begin{itemize}
\item ASC were most often predicted B cells;
\item MAIT cells were most often predicted \gd{} or CD8~T cells;
\item pDC were most often predicted cDC.
\end{itemize}
There were two exceptions: M6 predicted 96\% of pDC as HSPC, and M1 predicted ASC most often as ILC (0.269). Known-versus-removed AUROC ranged from \ctres{B}{M5}{unknown-auroc} (M5) to \ctres{B}{M4}{unknown-auroc} (M4). For M3, M5 and M6 it was below 0.5. AUROC is a rank probability: for M3 and M5, a removed-type cell received higher confidence than a known cell \emph{more often than not}; this does not mean that mean confidence was higher. M6's value (\ctres{B}{M6}{unknown-auroc}) is effectively chance. AUROC has no interval.

At M4's OP-err threshold, only \ctres{B}{M4}{unknown-false-accept-at-OPerr} of removed-type cells were accepted, but so were only \ctres{B}{M4}{coverage-at-OPerr} of known cells.

These are three near-neighbour immune populations. The result shows that confidence thresholds tuned for coverage do not flag them as unfamiliar. It does not show how these annotators behave for distant or genuinely novel cell types, which were not tested.

\paragraph{Arm~B scANVI changed sharply, and the change cannot be attributed.} Arm~B M6 performed markedly worse on known classes than Arm~A M6:
\begin{itemize}
\item accepted error at OP-cov rose from \ctres{A}{M6}{accepted-error-at-OPcov} to \ctres{B}{M6}{accepted-error-at-OPcov};
\item cross-lineage rate rose from \ctres{A}{M6}{cross-lineage-rate-at-OPcov} to \ctres{B}{M6}{cross-lineage-rate-at-OPcov};
\item RA accepted error rose from 0.220 to 0.516;
\item closed-set recall fell from 0.922 to 0.614 for CD4~T and from 0.731 to 0.572 for CD8~T;
\item ECE rose from 0.134 to 0.249;
\item natural-stratum precision for HSPC and ILC predictions was about 0.003.
\end{itemize}
The Arm~A M6 fit is a cached historical fit. The Arm~B fit was made in a different recovery session and adopted. With a single training seed per arm, and both reference training and the 50-epoch unlabelled query adaptation being stochastic, we cannot attribute this change to the class removal. Paired B:M6$-$B:M4 differences are valid within Arm~B but inherit the same caveat. Arm~B macro-F1 values are higher than Arm~A values for several methods (e.g.\ M4 \ctres{B}{M4}{macro-f1-closed} against \ctres{A}{M4}{macro-f1-closed}). These numbers average over different class sets (8 against 11), and the removed classes were among the hardest (ASC closed-set recall about 0.35 in Arm~A). The comparison therefore says nothing about whether removing classes helped.

\subsection{Natural unknowns: 20 erythroid cells, counts only}\label{sec:natunk}

In Arm~A, the natural stratum of the test studies contained 20 mapped cells outside $K$, all of class erythroid. This is one absent class from few donors. In 126 of the 1{,}000 bootstrap replicates there was no unknown cell at all, so its intervals are degenerate. We report counts only and draw no ranking or interval from them.

At OP-cov, the matched methods accepted the following numbers of these 20 cells with a known label (counts from the corrected run's false-acceptance fractions; most frequent predicted labels from the original run's unknown table):
\begin{itemize}
\item M1: 17 (mostly as platelet/MK);
\item M2: 13 and M3: 11 (mostly as CD4~T);
\item M4: 3 (most often predicted NK);
\item M5 (seed~0): 4. The predicted labels of these cells for the corrected fit are not in the supplied packet;
\item M6: 5 (most often predicted platelet/MK).
\end{itemize}
The M5 count has varied between fits on identical inputs: 6 in the original unseeded fit, 3 in the unseeded fit of the failed reproduction and 4 in the corrected seed-0 fit (Section~\ref{sec:m5}). An M4--M5 contrast on these cells is therefore not interpretable.

For P1 and P2 these numbers measure something else. Both released models have erythroid in their label space. The scoring code nevertheless counts any accepted prediction on these cells as a false acceptance, because ``unknown'' is defined against the matched Arm~A class set. In the original run's top-prediction distribution (P1 and P2 predictions are unchanged), P1 labelled 16 of the 20 cells erythroid and P2 labelled 18. P2's ``false acceptance'' of 20/20 and P1's of 14/20 at OP-cov therefore mostly count \emph{correct} erythroid calls. We exclude both from any comparison. Rescoring them with erythroid counted as correct would be a new analysis and was not performed. A secondary all-strata version of this analysis (126 cells) is again erythroid only and is not interpreted further.

\subsection{Calibration}

Table~\ref{tab:calib} (Appendix) lists Brier score and ECE. The comparison is meaningful only within the matched probabilistic methods. In Arm~A, ECE was 0.027 for M4, 0.081 for M5 (seed~0) and 0.134 for M6. Brier scores were M4 0.246, M5 0.258 and M6 0.299.

The other rows measure different things:
\begin{itemize}
\item \textbf{M3:} its row describes the reliability of kNN vote shares (ECE 0.065); vote shares are uncalibrated scores.
\item \textbf{M1 and M2:} their scores are uncalibrated by design.
\item \textbf{P1:} the Brier score uses CellTypist's one-vs-rest sigmoids renormalised to sum to one, while ECE uses the raw maximum sigmoid. P1's Brier (0.303) and ECE (0.121) are computed on different probability objects and cannot be compared with the matched rows.
\item \textbf{P2:} \texttt{coarser} predictions count as incorrect at their stated confidence.
\end{itemize}

\subsection{Platform panel: one donor per platform}

The full platform table is in \nolinkurl{platform.csv} of the scoring output. The corrected run's platform table was not in the packet supplied for this revision, so the values quoted here are original-run values. The table contains 1{,}000 sampled cells per platform dataset, of which 891--971 were scorable, and all come from a single donor. The original study labelled cells with a CellTypist and Seurat consensus \citep{desimone2025platforms}, so P1 agreement on this panel is partly circular. We report the most striking observations without causal explanation:
\begin{itemize}
\item P2 accepted 0.194 of cells on Parse Evercode v2 (closed-set accuracy 0.130) and 0.179 on ScaleBio (0.110).
\item M1's closed-set accuracy was 0.211 on Parse and 0.223 on ScaleBio. For the historical unseeded M5 fit it was 0.434 and 0.548. Between the two unseeded M5 fits, M5 labels changed for 1.2--3.0\% of the 1{,}000 cells on Parse, ScaleBio and PIP-seq, so these M5 values describe one fit only.
\item M2, M3 and M6 had closed-set accuracy between 0.92 and 0.98 on both of those platforms.
\end{itemize}
With one donor, these observations cannot separate platform effects from donor or sampling effects.

\subsection{Protein check (descriptive; replacement inputs)}

Label: \emph{descriptive; replacement inputs (amendment 2026-10-06-approved-protein-replacement)}.

The check used two author-processed totalVI files, 10k and 5k (Section~\ref{sec:amend}). Both passed every pre-specified (amendment) eligibility criterion. The six gates resolved a single protein class for:
\begin{itemize}
\item 6{,}375 of 6{,}855 cells (93.0\%) in the 10k file;
\item 3{,}762 of 3{,}994 cells (94.2\%) in the 5k file.
\end{itemize}
Table~\ref{tab:protein} reports results per file, without pooling and without intervals. The M5 rows use the corrected seed-0 fit's CITE predictions.

Two agreement measures are shown, and neither is an accuracy.
\begin{itemize}
\item \textbf{\nolinkurl{agreement_accepted}} compares fine RNA labels (14 classes) with six coarse protein classes by exact match. A MAIT call on a protein-CD8~T cell, or an ILC call on a protein-NK cell, counts as disagreement. This measure therefore mixes true disagreement with granularity mismatch.
\item \textbf{\nolinkurl{agreement_accepted_gateable_preds}} keeps only accepted predictions whose RNA label is one of the six gate names. Because it selects on the method's own output, methods that send difficult cells to non-gateable labels are not penalised.
\end{itemize}

\noindent The two measures have different denominators. \nolinkurl{agreement_accepted} is computed over all accepted, protein-resolved cells for that method (\nolinkurl{n_accepted_resolved}). \nolinkurl{agreement_accepted_gateable_preds} is computed over the subset of those cells whose accepted RNA label is one of the six gate names. That smaller, method-specific count is not reported in the outputs, and we did not reanalyse to obtain it.

For example, M4 had \nolinkurl{agreement_accepted} 0.904 on the 10k file over $n=6{,}009$ accepted, resolved cells, and \nolinkurl{agreement_accepted_gateable_preds} 0.978 over an unreported gateable subset of those cells. On the 5k file the values were 0.912 ($n=3{,}321$) and 0.950 (gateable subset, count not reported). M1 on the 5k file had 0.636 overall but 0.938 on gateable predictions, illustrating the granularity effect.

The accepted-and-resolved denominators differ between methods: from 5{,}027 (P2) to 6{,}331 (P1) cells on the 10k file, and from 3{,}306 (M1) to 3{,}726 (P1) on the 5k file. P2 again accepts every fine-level label ($\tau_\mathrm{cov}=-\infty$). Each method is therefore scored on a different subset of cells, and we draw no method ranking from this check.

Several further limits apply:
\begin{itemize}
\item Upstream author filtering may raise agreement.
\item Donor identity and overlap with the released models' training data are unknown.
\item About a third of the scTab gene universe, and about a fifth of $F$, was absent and set to zero, which is not neutral after z-scoring.
\item The CD16-mono protein class had only 46 (10k) and 43 (5k) cells.
\item The CLR row mean used 14 antibodies in the 10k file and 29 in the 5k file, so gate scales differ between files.
\end{itemize}

\begin{table}[t]
\centering
\caption{Protein check, per file. Descriptive; replacement inputs (amendment 2026-10-06-approved-protein-replacement). M5 rows are the corrected seed-0 fit. ``Accepted \& resolved'' is the per-method denominator of \nolinkurl{agreement_accepted} only. The gateable-prediction column uses a smaller, method-specific subset (accepted, resolved cells whose RNA label is one of the six gate names); that count is not recorded in the outputs and is not shown. Fine RNA labels are compared with six coarse protein classes by exact match, so neither column is an accuracy, and different methods are scored on different cells. No pooling, interval or ranking.}
\label{tab:protein}
\gentable{table_protein.tex}
\end{table}

\section{Discussion}

The question was how to choose an annotation workflow for a new blood cohort and when to accept a label. The measured results support four practical statements, each conditional on this reference, configuration, seed and set of four test studies.

\begin{enumerate}
\item \textbf{Among learned annotators trained on the same reference, the choice of method mattered little for closed-set agreement.} For logistic regression, kNN, CellTypist trained on our reference and scANVI, the paired intervals included zero; no difference was resolved at this precision, which falls short of showing equivalence. For M5 this held in the corrected seed-0 fit and in both earlier unseeded fits. Simpler rules (marker scores, centroids) agreed less with author labels. This is consistent with published reports that simple baselines remain competitive \citep{boiarsky2023,kedzierska2025,denadel2026}.
\item \textbf{For accepted error, the test study mattered more than the choice among the learned matched methods.} The same model's (M4) accepted error ranged from 0.021 to 0.179 across studies, while M3, M5 and M6 differed from M4 by 0.036 or less in pooled accepted error at OP-cov (point estimates of the results of record). This does not extend to the marker and centroid baselines, whose differences from M4 were larger, and the study-level comparison rests on four studies without an interval. A pooled number from a benchmark therefore tells an analyst little about their own cohort. The direction matches the published drop in agreement when models move to new studies \citep{cultrera2026hce}.
\item \textbf{A threshold tuned to accept 90\% of validation cells did not accept 90\% of test cells,} and a threshold tuned for 5\% validation error gave test errors from 0 to 16\% at very different coverages. Thresholds selected on validation data should be re-checked on each new study's own labelled subset if one exists. This interpretation follows from the transfer gaps we measured, but we did not test re-checking.
\item \textbf{At the coverage-oriented threshold, cells of missing near-neighbour types were mostly accepted with a confident, wrong, related label.} For M3 and M5, AUROC below 0.5 means a removed-type cell outranked a known cell in confidence more often than not; M6's AUROC (\ctres{B}{M6}{unknown-auroc}) is effectively chance. AUROC has no interval. This agrees qualitatively with earlier reports that posterior-based rejection can fail for missing populations \citep{abdelaal2019} and that closed-set annotators force absent types into existing labels \citep{hu2025matching}.
\end{enumerate}

For released models, the practical track shows what an analyst receives \emph{as distributed and mapped by our tables}. P2's 15.3\% \texttt{coarser} labels and P1's \gd{} label mapping illustrate how label vocabularies, independently of any classifier, shape apparent accuracy.

The M5 history adds a practical point about reproducibility. A seed stated in a protocol has to reach the estimator. Here the specified seed was never passed to the CellTypist solver, and the omission surfaced only when a reproduction compared per-cell outputs. Pooled M5 metrics moved by less than 0.01 between the two unseeded fits, but the per-cell labels, confidences and the small natural-unknown counts did not repeat.

\section{Limitations}\label{sec:limits}

\begin{itemize}
\item \textbf{Reference standard.} Accuracy means agreement with author labels mapped to the Cell Ontology; there was no adjudication. If author labels were themselves derived by reference mapping, agreement is inflated for methods that resemble that pipeline. The HIHA authors report CellTypist- and Seurat-guided annotation \citep{gong2025hiha}, and the platform panel used a CellTypist and Seurat consensus \citep{desimone2025platforms}. We did not verify how these labels were finally produced in the Census copies.
\item \textbf{Four fixed test studies.} Intervals are donor-bootstrap intervals within these studies, conditional on fixed thresholds and one training seed. They do not describe performance on new studies. Glaucoma and JDM contribute only 3 and 5 donors. Pooled values are weighted by cell count and dominated by HIHA and RA. No sensitivity analysis excluding HIHA (whose labels were partly guided by CellTypist) or reweighting studies was performed; only per-study values are reported.
\item \textbf{Bootstrap with absent classes.} Replicates with no unknown cell (126) give an undefined natural-unknown value, so that interval rests on 874 replicates and is conditional on an unknown cell being present (Table~\ref{tab:nan}). The only other undefined replicates are OP-err for M3 and M6, which was not attainable. A replicate missing a macro-F1 class scores that class 0 and is kept, which biases replicate macro-F1 downward by an unrecorded amount.
\item \textbf{Single training seed.} No training-seed variance is included for any method. This matters most for scANVI, whose Arm~A and Arm~B fits came from different sessions and differed sharply.
\item \textbf{M5 seed and convergence.} The original M5 fits were unseeded, which deviated from the protocol. The corrected fits use seed~0 and are the results of record. Their coefficients still depend on the seed. In the unseeded fits, at least one one-vs-rest sub-problem per fit stopped at \code{max\_iter=500} before meeting the tolerance, so the coefficients are a stopping point on a random-order SAG path. The iteration cap and solver are part of the specified method and were not changed. The three M5 fits (two unseeded, one seeded) are not a seed-variance estimate, and they do not isolate the seed as the only cause of the earlier disagreement (Section~\ref{sec:m5}).
\item \textbf{Pending reproduction of the corrected M5 results.} \Reprostatus. Even when a comparison report exists, it is an internal comparison by the same pipeline. No independent human scientific review or external replication has been carried out.
\item \textbf{Thresholds.} The thresholds were fixed on two validation studies whose mix includes the rare top-up stratum. M4's regularisation strength was chosen on the same studies, which is a mild double use of validation data that leaves the test data independent.
\item \textbf{Ties.} M3's tie handling makes its OP-cov threshold a 94\% (Arm~A) and 92\% (Arm~B) validation operating point, and makes its risk--coverage values depend on row order. Fixing this would require rescoring a primary metric, which needs approval.
\item \textbf{Natural unknowns.} These were 20 erythroid cells, so the study gives no evidence about natural novel-type detection. The M5 count changed between fits (6, 3 and 4 of 20).
\item \textbf{Simulated unknowns.} These cover three near-neighbour classes and were mostly top-up cells.
\item \textbf{Practical track.} Results mix training data, label vocabularies, gene spaces and our frozen label maps.
\item \textbf{Platform panel.} One donor per platform, with labels of uncertain provenance. The M5 platform values quoted are from the historical unseeded fit.
\item \textbf{Protein check.} Descriptive only: replacement files, author-filtered, zero-filled genes, donor identity and training overlap unknown, fine-versus-coarse label mismatch.
\item \textbf{Reference composition.} The reference includes in-vitro-activated PBMCs from the CVID study \citep{rodriguezubreva2022cvid}. The exact \texttt{disease == 'normal'} filter may omit some healthy cells, because the Census \texttt{disease} field can hold multiple values \citep{censusreleases}.
\item \textbf{Runtime and memory.} The protocol lists per-method runtime and peak memory, but this paper does not report them systematically (Section~\ref{sec:repro}).
\end{itemize}

\section{Reproducibility, data and code}\label{sec:repro}

\paragraph{Status.} The results of record are the corrected canonical run \runid, which combines seed-0 M5 fits with outputs of the other methods adopted unchanged from the original run \origrunid. \Reprostatus. The generated table below is filled in from the tiered comparison report when one is supplied to the asset generator; otherwise it states that none was supplied. Even when it is present, it compares the reproduction with the corrected canonical run within pre-set tolerances and is not independent external verification.

\begin{table}[H]
\centering
\caption{Tiered reproduction comparison (generated). Without a comparison report, it states that none was supplied.}
\label{tab:repro}
\gentable{table_reproduction.tex}
\end{table}

\paragraph{Commands.} From the repository root, the original full pipeline is:
\begin{verbatim}
make env          # uv sync --frozen (locked environment)
make data         # scripts/acquire_inputs.py acquire --root .
make test         # pipeline unit tests
CELLTRANSFER_BASELINE_MANIFEST=<baseline comparison manifest> \
  make reproduce  # scripts/reproduce.py --config configs/full.json
make paper        # scripts/build_paper.py (Tectonic, pinned bundle)
\end{verbatim}
\texttt{reproduce.py} runs, in order:
\begin{enumerate}
\item acquisition;
\item the data build;
\item Arms A and B fit and predict for M1--M6;
\item the CITE build (the amended totalVI builder);
\item scoring (\verb|score_all.py --reps 1000 --natural-unknown-scope natural|);
\item the protein check;
\item the comparison against the baseline manifest;
\item asset generation (\nolinkurl{make_paper_assets.py}) and the paper build.
\end{enumerate}
The corrected M5 stages use the repair driver recorded in the run manifest, with configuration \nolinkurl{configs/m5-seed-repair.json}:
\begin{verbatim}
uv run --frozen python scripts/m5_seed_repair.py --stage canonical \
  --config configs/m5-seed-repair.json --output <output dir>
uv run --frozen python scripts/m5_seed_repair.py --stage reproduction \
  --config configs/m5-seed-repair.json --output <output dir>
\end{verbatim}
For a run without a comparison report, the assets are generated with:
\begin{verbatim}
python scripts/make_paper_assets.py --score <score dir> \
  --protein <protein dir> --no-comparison \
  --eligibility <eligibility.json> --output paper
\end{verbatim}
After a comparison, \verb|--no-comparison| is replaced by \verb|--comparison <report.json>|, and the status statements in this paper switch automatically.

\paragraph{Full tables not reproduced here.} To keep tables readable, the large generated tables are not typeset. They are maintained as CSV files in the scoring output directory: \nolinkurl{per_study_test.csv} (per-study values), \nolinkurl{paired_differences_vs_M4.csv} (all paired differences), \nolinkurl{platform.csv} (platform panel), \nolinkurl{cell_counts.csv} (cell counts), \nolinkurl{risk_coverage_test.csv} and \nolinkurl{unknowns_allstrata_secondary.csv}. The corrected run's paired differences and cell counts are also in its \nolinkurl{results.json}.

\paragraph{Budgets.} The approved budgets are: local compute only; four CPU threads; a 12~GiB process-group memory watchdog; 7~GiB of study storage; at most two attempts per step; 8~h for recovery, 15~h for clean reproduction and 30~min for a paper-only rebuild. The M5 repair added no new runtime ledger. It is charged to the existing 54{,}000~s reproduction ledger with its single 30~s reserve (effective ceiling 53{,}970~s), and it permits exactly four seed-0 M5 fits (canonical A and B, reproduction A and B) with no fit retry. The corrected run's manifest records 10{,}490~s cumulative after it finished; this includes 7{,}085~s from the earlier reproduction lineage and both failed canonical attempts.

\paragraph{Recorded resource use.} The original run's \texttt{/usr/bin/time} records cover only the steps it recomputed, and they measure only the direct child process, so they understate aggregate memory under parallelism:
\begin{itemize}
\item scoring: 557~s, with a maximum resident set of about 0.57~GB;
\item CITE build: 16.2~s, about 5.3~GB.
\end{itemize}
Arm~A fit and predict timings are in the historical run records. Arm~B M4--M6 wall times were recorded in the adopted session: the historical unseeded M5 fit took 564~s and the M6 fit 713~s. In the corrected run, the Arm~B M5 fit took 501~s, scoring 613~s and the protein check 11~s (wall times from its receipt). Each is a single run, so we make no runtime comparison.

\paragraph{Failures and attempts (all records kept).}
\begin{itemize}
\item An early data build was aborted, and a second stopped on a disk-guard timeout before any query data were pulled. The full build D05 used the same code and seed, with resumable checkpoints that were tested for equivalence.
\item The first Arm~B M4 fit failed for lack of disk space. A pipeline test later launched one unintended Arm~B M4 fit. It was killed on detection without producing a model, and 60~s were charged to the budget. The adopted Arm~B M4 fit is attempt~2.
\item The only pre-specified 10x CITE-seq file that was requested returned HTTP 403. The replacement CITE build was attempt~2, the last one permitted.
\item Pan-Human Azimuth and scGPT were not run.
\item The reproduction \failrunid{} failed its comparison gate because of M5 (Section~\ref{sec:m5}).
\item Two corrected canonical attempts failed before run \runid{} completed (Section~\ref{sec:m5}). Both are retained and their time is charged.
\end{itemize}

\paragraph{Code and data availability.}
\begin{itemize}
\item \textbf{Code} is in this repository: \texttt{companion/} holds the methods, scoring and protein check, and \texttt{scripts/} holds orchestration, the M5 repair driver, the asset generator and the paper build.
\item \textbf{Code hashes} (sha256) used for the corrected run:
\begin{itemize}
\item \nolinkurl{score_all.py}: \nolinkurl{20f457204050bdac2c0d99596c43b9bb3214e9ff974ba24041f50eeab4e54d9f}
\item \texttt{evaluate.py}: \nolinkurl{e5fae4e8646821887d4555bc7746ba1bd1b4a259d172197590808d7d215a59a2}
\item \nolinkurl{protein_check.py}: \nolinkurl{9458409448615ebce2bab0b189ef113c57371d57af647ef27c2c1c2c7442ba59}
\item \texttt{methods.py} with the seed argument: \nolinkurl{d2cb0cdef8a25cdc98df43ffdff3c9fcf6040972cab51b8ad5416be7da2f167e}; original, unseeded: \nolinkurl{914989c507d02ffe04886907b0812fd3c80ac597c8c87a657890923f35cf3c14} (preserved as \nolinkurl{evidence/reviews/m5-original-methods-1e95aac.py})
\end{itemize}
\item \textbf{Census data} are CC BY 4.0 and are accessed by dataset ID from LTS 2025-11-08 \citep{czi2025cellxgene,censusreleases}.
\item \textbf{Released models:} CellTypist \nolinkurl{Immune_All_Low} v2 \citep{celltypistmodels} and the scTab run5 checkpoint \citep{sctabrepo}.
\item \textbf{Aggregate results} of the corrected run are in \nolinkurl{evidence/reviews/m5-canonical-448f32386dfe45588fefa84c5d486703/}. Its \nolinkurl{results.json} there is byte-identical to the run output; the manifest and receipt there have absolute path prefixes replaced, with the original hashes recorded in \nolinkurl{bundle-index.json}. The original run's aggregate results remain in \nolinkurl{evidence/results-29879296/}. Per-cell protein classes and CITE predictions are not published (Section~\ref{sec:amend}).
\end{itemize}

\section{M5 seed correction and versioning}\label{sec:m5}

This section records how the M5 results of record changed and what remains unresolved. The corrected results carry the version label \code{m5-seed0-v2}. The deviation is logged in \nolinkurl{evidence/protocol-deviations-m5-seed0-v2.md}; superseded claims are mapped in \nolinkurl{evidence/claims-history/m5-seed0-v2-supersession.json}.

\paragraph{The original M5 fits were unseeded.} The protocol states ``Seeds: 0 for model training''. The original M5 code called \code{celltypist.train} without a seed. CellTypist selects the SAG solver for more than 50{,}000 cells and passes \code{random\_state} to scikit-learn only when it is given, so the saved original estimators record \code{random\_state=None}, with \code{max\_iter=500}, \code{tol=1e-4}, \code{C=1.0} and \code{n\_jobs=4} (failure review, finding E4). M2/M3 (PCA), M4 and M6 were already seeded with 0. An earlier version of this paper stated that the training seed was 0 for all methods; that was incorrect for M5.

\paragraph{The reproduction gate failed on M5.} A reproduction lineage made fresh fits of M1--M6; its final run \failrunid{} carried the M5 fits forward from parent run \nolinkurl{727c2514a006421286527ce1d2f22148} by hash and was compared with the original run \origrunid{} under the pre-specified tolerances (\nolinkurl{protocol/tolerances-v2.json}). According to the retained comparison summary (\nolinkurl{evidence/reviews/m5-failure-78542e9d412548edbfe42b38d452f7a8/comparison-failure-summary.json}, a derivative of report sha256 \nolinkurl{2a663866746ee682b93cadcfaebba56715dbb6cf308bc78bdd86301f689d6dc8}), the comparison ran 1{,}523 checks and reported overall status ``partially reproduced'', with 54 gated breaches and 0 unsupported checks. All 54 breaches were M5:
\begin{itemize}
\item 17 per-file label checks (8 Arm~A files and 9 Arm~B files exceeded the allowed number of disagreeing cells; the other M5 label files stayed within it);
\item 32 per-file confidence checks (all 17 Arm~A and all 15 Arm~B M5 files);
\item 5 metric checks, all Arm~A natural-unknown quantities: false acceptance at OP-cov, its upper interval bound, the AUROC, and the paired mean difference and upper bound against M4.
\end{itemize}
M1--M4, M6, P1 and P2 passed their declared gated checks. The comparison did not exit~0, so the gate failed. The inputs to both M5 fits (training matrix, labels, features and scaler) were identical, and the coefficients differed; the divergence therefore entered at the solver fit (failure review, E3). In the comparison outputs, the other pooled M5 metrics moved by less than 0.01 between the two unseeded fits, and the paired M5$-$M4 macro-F1 interval included zero in both.

\paragraph{The approved correction.} The user approved one scientific change on 7 October 2026: pass \code{random\_state=0} to the M5 solver. The solver, \code{max\_iter=500}, data, features, labels, thresholds, metrics, bootstrap, tolerances and resource limits were unchanged. The approval allows four seed-0 M5 fits (canonical A and B, reproduction A and B), two scoring stages, no fit retry and no new runtime ledger. All historical and failed outputs are retained. The corrected M5 results replace the original M5 results as results of record; this is a material change to the M5 result of record and was approved as such.

\paragraph{Canonical attempts.} Three canonical attempts were made, and all are retained and charged.
\begin{enumerate}
\item Attempt \nolinkurl{c098a3420d0a4c99b6d2b8d65c667af7} stopped after 303~s, before any M5 fit was invoked. The recorded recovery approval changed only engineering inputs: the M6 checksum-sidecar paths were moved, with identical bytes, so that they no longer matched the rule that excludes M5 artefacts from adoption.
\item Attempt \nolinkurl{3f18e178b78d4d75b595174cfa00cf99} completed the seed-0 Arm~A fit and then stopped on a checker-encoding mismatch. The historical scaler digests had been computed over dtype, shape and bytes, while the checker first hashed raw bytes only. A read-only comparison showed that the genes, classes and scalers were identical to the historical ones. The checker's digest encoding was corrected, with no change to model, data, parameters or tolerances, under a recorded user approval.
\item Run \runid{} continued from attempt~2. It retained the completed Arm~A fit (\verb|executed: false|, source step hash recorded), ran only the Arm~B fit, and then ran the predictions, scoring and protein check (Section~\ref{sec:provenance}). It completed with exit code~0.
\end{enumerate}
No additional fit was made and no scientific setting changed across these attempts. Two of the four permitted M5 fits have now been used, one in attempt~2 and one in run \runid.

\paragraph{Descriptive comparison of the three M5 fits.} Table~\ref{tab:m5hist} sets the two historical unseeded fits next to the corrected fit. The historical columns are copied from the failure review (E1, E2); the corrected column is generated. This is descriptive only, and three fits do not estimate seed variance.

\begin{table}[H]
\centering
\small
\caption{M5 across three fits on identical inputs (Arm~A unless stated). Historical columns: original unseeded fit (run \origrunid) and unseeded fit from the failed reproduction (\failrunid), copied from the failure review. Corrected column: seed-0 fit, results of record, generated from the build outputs. Descriptive; no interval except the paired difference; not a seed-variance estimate.}
\label{tab:m5hist}
\begin{adjustbox}{max width=\textwidth}
\begin{tabular}{lccc}
\toprule
Quantity & Original (unseeded) & Reproduction (unseeded) & Corrected (seed 0) \\
\midrule
Macro-F1, closed & 0.688 & 0.689 & \ctres{A}{M5}{macro-f1-closed} \\
Coverage at OP-cov & 0.798 & 0.798 & \ctres{A}{M5}{coverage-at-OPcov} \\
Accepted error at OP-cov & 0.089 & 0.089 & \ctres{A}{M5}{accepted-error-at-OPcov} \\
AURC & 0.0526 & 0.0525 & \ctres{A}{M5}{aurc} \\
Natural-unknown false acceptance (of 20) & 0.30 (6) & 0.15 (3) & \ctres{A}{M5}{unknown-false-accept-at-OPcov} \\
Natural-unknown AUROC & 0.803 & 0.826 & \ctres{A}{M5}{unknown-auroc} \\
Arm~B macro-F1, closed & 0.747 & 0.748 & \ctres{B}{M5}{macro-f1-closed} \\
Arm~B removed-type false acceptance & 0.834 & 0.834 & \ctres{B}{M5}{unknown-false-accept-at-OPcov} \\
Paired M5$-$M4 macro-F1 & $-0.006$ [$-0.014$, 0.002] & $-0.005$ [$-0.014$, 0.003] & \ctres{A}{M5}{diff-vs-A-M4-macro-f1-closed} \ctci{A}{M5}{diff-vs-A-M4-macro-f1-closed} \\
\bottomrule
\end{tabular}
\end{adjustbox}
\end{table}

\paragraph{What remains unresolved.}
\begin{itemize}
\item The unseeded solver is the leading and sufficient candidate explanation for the earlier disagreement, but it is not proven to be the only cause. Other contributors, such as \code{n\_jobs=4} one-vs-rest parallelism, BLAS threading or environment differences in the historical fit, were not excluded. A synthetic engineering fixture with SAG and \code{max\_iter=500} gave identical outputs under three ambient NumPy seeds once the estimator seed was set, but that is not a test on study data.
\item Whether seeding makes the real M5 fits repeatable within the tolerances is tested only by the seed-controlled reproduction. That reproduction is planned to refit M5 in both arms in a separate checkout, with no copy of any canonical M5 artefact, and to adopt the other methods' completed fits from the failed reproduction lineage (\nolinkurl{727c2514a006421286527ce1d2f22148} / \failrunid), which passed their gated checks against the original run. Those adopted fits are not new fits made for this correction. \Reprostatus.
\item If that comparison does not pass, the approved plan stops without further seeds, fits, solver or iteration changes, or tolerance changes, and the outcome will be reported as it is.
\end{itemize}

\section{Protocol amendments and required disclosures}\label{sec:amend}

\paragraph{M5 training seed (m5-seed0-v2).} Approved correction of the M5 result of record (Section~\ref{sec:m5}; \nolinkurl{evidence/protocol-deviations-m5-seed0-v2.md}).

\paragraph{Natural-unknown scoring (D-1a).} Under the approved resumption plan, Arm~A natural-unknown metrics use natural-stratum cells only, as the frozen protocol's ``natural stratum only'' rule requires.

\paragraph{Protein replacement amendment.} The following disclosures are required by the amendment and are reproduced as worded there, with the realised values filled in.

\begin{enumerate}
\item The study's protocol (v1.0, frozen locally on 2026-10-05; not externally registered) pre-specified four 10x CITE-seq files. Only the first was requested; it returned HTTP 403 on 2026-10-06. The other three were not requested, so their availability is unknown. Under an amendment approved on 6 October 2026 (amendment file sha256 \nolinkurl{e5ba0c347b66cba1c76dd249e386789c26e1916a95bc9c82a0208a696c89eae1}), we instead used author-processed files from the totalVI reproducibility repository (\url{https://github.com/YosefLab/totalVI_reproducibility}, commit \nolinkurl{27d65a9c49182f39df0c4325249bc7dc8e663dd0}), sha256:
\begin{itemize}
\item \nolinkurl{pbmc_10k_protein_v3.h5ad}: \nolinkurl{5f08b8575febf9e04b209b94eb43f6335f1e33b0985cdcf44adf7320f6243c69}
\item \nolinkurl{pbmc_5k_protein_v3.h5ad}: \nolinkurl{a1bf51e070d24b39627ea4de9b3e489a4637ff795e1061e0733e26c3baec8847}
\end{itemize}
Both files were used from previously downloaded bytes that matched these pins (recorded as \texttt{cached}), as the amendment permits.
\item The 5k file is 10x \nolinkurl{5k_pbmc_protein_v3}, not the pre-specified \nolinkurl{5k_pbmc_protein_v3_nextgem}. No replacement was used for \nolinkurl{10k_PBMCs_TotalSeq_B_3p} or the 5$'$ \nolinkurl{vdj_v1_hs_pbmc3}.
\item Both files were filtered by the totalVI authors (doublet removal, mitochondrial and complexity caps, manually chosen ADT library-size limits, isotype-control removal, gene filtering). These choices cannot be undone. They may raise the agreement compared with unfiltered data.
\item Genes were mapped from symbols to Ensembl IDs using only the Census 2025-11-08 table. The fraction of matched-track features present was 79.1\% (10k) and 77.2\% (5k) for $F_A$, and 79.2\% and 77.8\% for $F_B$. The fraction of scTab genes present was 66.9\% (10k) and 66.3\% (5k). Absent genes were set to zero, which is not neutral after z-scoring. Differences between methods on these files therefore mix method behaviour with robustness to missing genes.
\item Donor identity is unknown for both files, so we cannot say whether they come from independent donors. Training overlap with the released models is unknown. Results are descriptive, per file, with no interval and no generalisation or holdout claim. Protein gates are a coarse orthogonal check, not ground truth.
\item Every CITE result is labelled ``descriptive; replacement inputs (amendment 2026-10-06-approved-protein-replacement)''. Both files were eligible.
\end{enumerate}
\noindent Licence: totalVI reproducibility repository: no licence file found (GitHub API, 2026-10-06); terms of the 10x source data not verified by us. The h5ad files, derived query files, ADT tables, per-cell protein classes and per-cell CITE predictions are therefore not redistributed.

\bibliographystyle{unsrtnat}
\bibliography{references}

\clearpage
\appendix
\section{Additional generated tables}

These tables are generated directly from the score outputs the paper is built from; M5 rows are the corrected seed-0 fits. The per-study, paired-difference, platform and cell-count tables are not typeset because they are too large to read at page width; they are maintained as CSV files (Section~\ref{sec:repro}). The same conventions apply throughout. Intervals are 95\% donor-bootstrap percentile intervals within the four test studies, conditional on fixed thresholds and one training seed. ``Not attainable'' follows the protocol's OP-err rule. Arm~A and practical-track unknown-false-acceptance entries refer to 20 erythroid cells and must not be compared (Section~\ref{sec:natunk}).

\begin{table}[H]
\centering
\caption{Validation thresholds. P2's $\tau_\mathrm{cov}$ is $-\infty$ because its eligible validation coverage cap (0.798) is below 0.90. M3's validation coverage at $\tau_\mathrm{cov}$ exceeds 0.90 because of ties.}
\label{tab:thresholds}
\gentable{table_thresholds.tex}
\end{table}

\begin{table}[H]
\centering
\caption{Calibration. M1 and M2 scores are uncalibrated by design. M3 is vote-share reliability. P1's Brier score and ECE are computed on different probability objects (renormalised versus raw sigmoid). P2 counts \texttt{coarser} predictions as incorrect.}
\label{tab:calib}
\gentable{table_calibration.tex}
\end{table}

\begin{table}[H]
\centering
\caption{Practical-track label outcomes after mapping by the frozen label tables.}
\label{tab:practical}
\gentable{table_practical_labels.tex}
\end{table}

\begin{table}[H]
\centering
\caption{Bootstrap replicates with an undefined value. A count of 1{,}000 for M3 and M6 at OP-err means the operating point was not attainable. A count of 126 for Arm~A and practical unknown false acceptance means those replicates contained no erythroid cell.}
\label{tab:nan}
\gentable{table_bootstrap_nan.tex}
\end{table}

\end{document}
